% Options for packages loaded elsewhere
\PassOptionsToPackage{unicode}{hyperref}
\PassOptionsToPackage{hyphens}{url}
\documentclass[
]{article}
\usepackage{xcolor}
\usepackage{amsmath,amssymb}
\setcounter{secnumdepth}{-\maxdimen} % remove section numbering
\usepackage{iftex}
\ifPDFTeX
  \usepackage[T1]{fontenc}
  \usepackage[utf8]{inputenc}
  \usepackage{textcomp} % provide euro and other symbols
\else % if luatex or xetex
  \usepackage{unicode-math} % this also loads fontspec
  \defaultfontfeatures{Scale=MatchLowercase}
  \defaultfontfeatures[\rmfamily]{Ligatures=TeX,Scale=1}
\fi
\usepackage{lmodern}
\ifPDFTeX\else
  % xetex/luatex font selection
\fi
% Use upquote if available, for straight quotes in verbatim environments
\IfFileExists{upquote.sty}{\usepackage{upquote}}{}
\IfFileExists{microtype.sty}{% use microtype if available
  \usepackage[]{microtype}
  \UseMicrotypeSet[protrusion]{basicmath} % disable protrusion for tt fonts
}{}
\makeatletter
\@ifundefined{KOMAClassName}{% if non-KOMA class
  \IfFileExists{parskip.sty}{%
    \usepackage{parskip}
  }{% else
    \setlength{\parindent}{0pt}
    \setlength{\parskip}{6pt plus 2pt minus 1pt}}
}{% if KOMA class
  \KOMAoptions{parskip=half}}
\makeatother
\setlength{\emergencystretch}{3em} % prevent overfull lines
\providecommand{\tightlist}{%
  \setlength{\itemsep}{0pt}\setlength{\parskip}{0pt}}
\usepackage{bookmark}
\IfFileExists{xurl.sty}{\usepackage{xurl}}{} % add URL line breaks if available
\urlstyle{same}
\hypersetup{
  pdftitle={Posts as Events: Intraday Cryptocurrency Volatility Around a Head of State's Social-Media Communications},
  pdfauthor={{[}Author name and affiliation --- to complete{]}},
  hidelinks,
  pdfcreator={LaTeX via pandoc}}

\title{Posts as Events: Intraday Cryptocurrency Volatility Around a Head
of State's Social-Media Communications}
\author{{[}Author name and affiliation --- to complete{]}}
\date{June 2026}

\begin{document}
\maketitle

\subsection{Abstract}\label{abstract}

We study whether the social-media communications of a sitting head of
state carry short-horizon, market-relevant information for
cryptocurrency markets. Using the full corpus of 33,712 Truth Social
posts by Donald J. Trump (February 2022 -- June 2026) and hourly prices
for Bitcoin, Ethereum and Solana, we construct a text-only event signal
--- combining a financial-domain sentiment model, a large-language-model
market-impact stance read, a keyword topic taxonomy and a causal novelty
score --- and conduct a burst-level event study. We find that
market-relevant posts are followed by a modest but robust increase in
1--4 hour \textbf{realised volatility} (Bitcoin: +5.2 bp over four
hours, day-clustered regression \emph{p} = 0.007), which survives a null
matched on the prevailing volatility regime and, critically, is
\textbf{absent for the same author's non-market-relevant posts} (a
content placebo run on five times as many events). The effect does
\textbf{not} extend to 24 hours (that horizon is consistent with
ordinary volatility clustering) and does \textbf{not} translate into
predictable return \emph{direction}: neither the directional signal nor
any post-event trading rule produces a tradeable edge net of transaction
costs, a null that is robust to using a sentiment model whose training
predates the events. A small class of explicit ``market-directive''
posts --- including the widely reported ``THIS IS A GREAT TIME TO BUY''
message of 9 April 2025 --- preceded large rallies, but with only two
such events these are reported as a qualitative case study, not as
evidence of predictability. The contribution is methodological as much
as empirical: an intraday, full-corpus, text-only, pre-registered design
with an endogenous-timing control and a content placebo, applied to a
setting (Truth Social → crypto, the 2025--26 trade-policy regime) that
prior daily, hand-curated event studies do not resolve.

\textbf{Keywords:} event study; realised volatility; social media;
cryptocurrency; political communication; volatility clustering.

\begin{center}\rule{0.5\linewidth}{0.5pt}\end{center}

\subsection{1. Introduction}\label{introduction}

A distinctive feature of the 2025--26 market environment is that a
single, high-authority actor --- the President of the United States ---
routinely posts unscheduled statements on a social-media platform (Truth
Social), some of which concern, or directly precede, market-moving
policy. Three episodes drew particular public attention. On 3 March 2025
a post announcing a ``Strategic National Crypto Reserve'' was followed
by a sharp Bitcoin rally. On 9 April 2025, roughly four hours before the
announcement of a 90-day pause in reciprocal tariffs, the President
posted ``THIS IS A GREAT TIME TO BUY!!! DJT''; equity and crypto markets
rallied sharply that afternoon, and several legislators publicly called
for insider-trading and market-manipulation investigations. On 10
October 2025 a post announcing 100\% tariffs on China was followed by a
\textasciitilde12\% intraday fall in Bitcoin and the largest single-day
liquidation event in the asset's history to that point.

These episodes raise a measurement question that this paper addresses
quantitatively and neutrally: \textbf{do this actor's posts carry
short-horizon, market-relevant information for cryptocurrency, and if
so, in what form --- direction, volatility, or neither?} We make no
claim about intent; we report the public allegations around specific
posts as dated facts and study the \emph{abnormal returns and
volatility} around posts, not the question of manipulation.

We answer the measurement question with a burst-level event study on the
full post corpus, joined to hourly crypto prices and a daily macro
panel. Three pre-registered research questions structure the analysis:

\begin{itemize}
\tightlist
\item
  \textbf{RQ1 (volatility).} Do market-relevant post bursts predict
  elevated 1--4 hour forward realised volatility \emph{beyond} what the
  prevailing (trailing) volatility already implies?
\item
  \textbf{RQ2 (direction).} Does a text-derived directional signal
  predict the \emph{sign} of forward returns?
\item
  \textbf{RQ3 (tradeable edge).} Does a post-event trading rule beat
  zero net of realistic transaction costs out of sample?
\end{itemize}

Our answers are: \textbf{yes} (modest, robust, content-specific) to RQ1;
\textbf{no} to RQ2 and RQ3. The most defensible single sentence is that
\emph{market-relevant posts move how much crypto moves over the next few
hours, not which way, and not by enough to trade after costs.}

\textbf{Contributions.} Relative to the closest prior work (Section 2),
which conducts \emph{daily} event studies on \emph{hand-curated}
policy-event lists, we contribute: (i) an \textbf{intraday}
(1/4/24-hour) design that resolves an effect invisible at daily
resolution; (ii) a \textbf{systematic, text-only classification of the
entire 33k-post corpus} rather than a curated event list, removing
event-selection discretion; (iii) an explicit \textbf{endogenous-timing
control} --- matching the comparison set on prevailing volatility,
because a high-profile actor may post precisely when markets are already
moving; and (iv) a \textbf{content placebo} that isolates whether the
effect is about market \emph{content} or merely about the act of
posting. We frame the study as exploratory/observational; the data --- a
single dominant policy regime, two directive-class events, and crypto
assets that co-move at \textasciitilde0.92 --- cannot bear
strong-confirmatory claims, and we say so throughout.

\begin{center}\rule{0.5\linewidth}{0.5pt}\end{center}

\subsection{2. Related work}\label{related-work}

\textbf{Trump communications and market volatility.} A mature literature
studies the president's \emph{Twitter} output and \emph{equity/bond}
markets. Nishimura, Dong \& Sun (2021, \emph{Journal of Financial
Research}) study the second-moment response of US equities to Trump
tweets using daily realised volatility decomposed into continuous and
jump components, finding that tweeting raises volatility and jump
tail-risk in the later sample. The practitioner ``Volfefe Index''
(JPMorgan, 2019) linked tweets to intraday Treasury-rate volatility; its
citable academic anchor, Klaus \& Koser (2021, \emph{Finance Research
Letters}), finds the index adds predictive power for European equity
returns, heterogeneously and non-linearly. Born, Myers \& Clark (2017,
\emph{Algorithmic Finance}) run a daily firm-level event study on 15
tweets, finding transitory, noise-trader-driven abnormal returns. On the
trade-war channel specifically, Burggraf, Fendel \& Huynh (2020,
\emph{Applied Economics Letters}) find US--China tweets negatively
predict the S\&P 500 and positively predict the VIX. Two studies move
toward our intraday, topic-classified design: Gjerstad et al.~(2021,
\emph{Decision Support Systems}) use precise tweet timestamps with
high-frequency equity data and LDA topics, finding the market falls and
uncertainty rises after tweets, strongest for the ``trade war'' topic;
and Perico Ortiz (2023, \emph{Journal of Economics and Finance}) runs an
intraday VIX event study (5-minute bars, to five hours post-tweet) with
unsupervised topic clustering, finding foreign-policy/trade tweets raise
uncertainty most. This literature is \emph{Twitter + equities/bonds,
first term}; we are \emph{Truth Social + crypto, 2022--26}.

\textbf{Trump and crypto.} The closest direct work is a cluster of SSRN
working papers by Krause (2025) --- e.g., \emph{Trump-Era Cryptocurrency
Policy Events and Bitcoin Returns}, \emph{Trump, Tokens, and Tailwinds},
\emph{Tariffs, Tokens, and Turmoil} --- which conduct \textbf{daily}
event studies on \textbf{hand-curated} policy/advocacy events, with
multi-day cumulative-abnormal-return windows. We differ on four axes
(intraday resolution; full-corpus systematic classification rather than
curated events; a focus on the second moment with an explicit treatment
of volatility \emph{clustering}; and an endogenous-timing-matched null).
We therefore do not claim to be the first study of Trump and crypto; our
novelty is the \emph{joint} combination of these design choices in this
regime.

\textbf{Social-media sentiment and crypto.} A broad literature relates
aggregate Twitter sentiment to crypto returns, volatility and liquidity,
and documents that single high-profile actors (e.g., Elon Musk) generate
abnormal returns in named coins. This work typically uses aggregate
retail sentiment rather than a single high-authority actor whose posts
can \emph{precede the policy they concern} --- the feature that
motivates our directive case study.

\begin{center}\rule{0.5\linewidth}{0.5pt}\end{center}

\subsection{3. Data}\label{data}

\textbf{Posts.} We use a pinned snapshot (33,712 posts, 14 Feb 2022 -- 2
Jun 2026; SHA recorded in the project manifest) of a public archive of
@realDonaldTrump's Truth Social posts. Each post carries an ISO-8601 UTC
timestamp, HTML content, engagement counts, and media flags. Because the
upstream archive refreshes every few minutes, all analysis runs against
the frozen snapshot for reproducibility.

\textbf{Prices.} Hourly OHLCV for BTCUSDT, ETHUSDT and SOLUSDT
(Binance), spanning the post period. Bitcoin history is drawn from a
local parquet lake; Ethereum and Solana, which lack deep history in the
lake, are fetched from the Binance public REST API (same venue, so the
series are consistent by construction). A daily macro panel (DXY, gold,
S\&P 500, Nasdaq-100, US 10-year yield, VIX, and the three crypto
assets) covers December 2024 onward.

\textbf{Event-time joins and no look-ahead.} Posts are joined to prices
strictly by their post timestamp floored to the hour; all forward
returns and volatilities are measured from that hour's \emph{close}
onward (a conservative choice that misses the first within-hour minutes
of any reaction). We validate the join against a ground-truth event ---
the ``Liberation Day'' tariff sell-off of 2--3 April 2025, in which
Bitcoin fell \textasciitilde2.4\% in a single evening hour ---
confirming that timestamps align and that no timezone error contaminates
the windows.

\begin{center}\rule{0.5\linewidth}{0.5pt}\end{center}

\subsection{4. Signal construction (text
only)}\label{signal-construction-text-only}

All independent variables are computed from post text and timestamp
\textbf{only}; no realised return ever enters a label. This is
essential: a label informed by the outcome it is meant to predict would
make the analysis circular.

\begin{itemize}
\tightlist
\item
  \textbf{Topic taxonomy.} A frozen keyword taxonomy assigns each post
  to economic topics (tariffs/trade, China, Fed/rates, crypto, dollar,
  energy, markets, fiscal), geopolitics, and two content classes
  motivated by the 2025 episodes: \emph{market-directive} (explicit
  buy/sell or market predictions, e.g.~``great time to buy'') and
  \emph{reassurance} (``be cool'', ``don't panic''). A post is
  \emph{market-relevant} if it hits any economic topic (≈8\% of posts).
\item
  \textbf{Financial sentiment.} FinBERT (ProsusAI/finbert) yields class
  probabilities; \texttt{finbert\_score\ =\ P(pos)\ −\ P(neg)}. On
  market-relevant posts this averages −0.01, versus +0.37 for a generic
  lexical sentiment model (VADER) --- i.e.~the domain model removes a
  strong positive bias driven by the author's superlative style.
\item
  \textbf{LLM market-impact stance.} For the \textasciitilde1,851
  market-relevant, non-boilerplate posts (the cases where idiom
  matters), an instruction-tuned LLM labels the \emph{direction a trader
  would lean for risk assets} on reading the post (−2\ldots+2), a
  conviction (0--1), and flags for policy-signal and market-directive
  content. We emphasise this is a market-impact read from the text, not
  hindsight.
\item
  \textbf{Novelty.} A causal 7-day novelty score (1 − maximum cosine
  similarity to the embeddings of posts in the trailing week)
  distinguishes first-of-kind statements from repetitive ones.
\item
  \textbf{Bursts.} Because posts arrive in clusters (ten posts in seven
  minutes is common), the unit of observation is the \emph{burst} ---
  posts less than 30 minutes apart --- anchored to the first post's
  hour. Within-burst moves cannot be attributed to one post.
\end{itemize}

We note two disclosed limitations of the signal. First, the directional
\texttt{signal} blends an LLM market-impact read (where available) with
FinBERT tone (otherwise); we therefore also report results using a
\textbf{FinBERT-only} signal, whose training predates these events and
is thus immune to any leakage through an LLM's parametric memory of how
famous posts played out. Second, engagement counts are \emph{not} known
at post time and are excluded from all predictive specifications.

\begin{center}\rule{0.5\linewidth}{0.5pt}\end{center}

\subsection{5. Method}\label{method}

The analysis was \textbf{pre-registered} before the signal-dependent
results were computed (with one dated amendment that reframed the study
as exploratory/observational and adopted the robustness checks below);
the pre-registration timestamp is the evidence of pre-commitment.

\textbf{Targets.} For horizons \emph{k} ∈ \{1, 4, 24\} hours we compute
forward log return and forward realised volatility (the root of summed
squared hourly log returns over \emph{k}), plus a trailing-24-hour
realised volatility \emph{known at event time}.

\textbf{Primary inference (RQ1).} A cluster-robust OLS regression on all
hourly bars of forward volatility on an event dummy, controlling for
trailing volatility and hour-of-day, with standard errors clustered by
calendar day. The trailing-vol control is the key identification device:
a high-profile actor may post \emph{because} volatility is already
elevated, so the question is whether forward volatility rises
\emph{beyond} what trailing volatility predicts.

\textbf{Corroborating nulls.} Two bootstrap nulls (2,000 draws): one
matching the comparison hours on hour-of-day \emph{and} trailing-vol
quintile (an i.i.d. null), and one resampling whole calendar-day blocks
to reproduce the event sample's day-clustering. We report a balance
table (event vs.~pool trailing volatility) with every volatility claim.

\textbf{RQ2/RQ3.} For direction, cluster-robust regressions of forward
returns on the signal, novelty and topic dummies, with
Benjamini--Hochberg false-discovery correction across the test grid. For
tradeable edge, a sign-of-signal rule on a final-20\%-by-time
out-of-sample holdout, held four hours, net of a 0.24\% round-trip cost
(0.1\% taker × 2 + 2 bp slippage × 2), judged by a block-bootstrap 95\%
Sharpe confidence interval whose lower bound must exceed zero. RQ3 is
reported for both the LLM-blended and FinBERT-only signals, and split by
policy-signal status.

\begin{center}\rule{0.5\linewidth}{0.5pt}\end{center}

\subsection{6. Results}\label{results}

Reported \emph{p}-values are \textbf{nominal} (uncorrected for multiple
comparisons) except in RQ2/RQ3, where a Benjamini--Hochberg
false-discovery correction is applied. For the RQ1 volatility result the
multiplicity safeguard is not a corrected threshold but the
\emph{consistency} of the effect across two horizons (1 h, 4 h) × three
assets × two independent nulls, together with the content-placebo
contrast (§6.2); we read the nominal values in this exploratory/
observational spirit rather than as confirmatory significance tests.

\subsubsection{6.1 RQ1 --- a robust 1--4 hour volatility
effect}\label{rq1-a-robust-14-hour-volatility-effect}

Market-relevant post bursts are followed by elevated short-horizon
realised volatility. In the cluster-robust regression (the primary
inference), the Bitcoin event dummy is \textbf{+3.06 bp at 1 hour
(\emph{p} = 0.018)} and \textbf{+5.18 bp at 4 hours (\emph{p} = 0.007)}
over the full sample, and +3.37 bp (\emph{p} = 0.041) / +6.33 bp
(\emph{p} = 0.015) in the in-office (2025+) sub-sample. The
\textbf{24-hour} dummy is insignificant and near zero (−0.89 bp,
\emph{p} = 0.795), consistent with the 24-hour window capturing ordinary
volatility clustering rather than a post effect. The two bootstrap nulls
corroborate the 1--4 hour result (event vs.~matched null: 4-hour +87.6
bp vs.~+82.8 bp). The balance table shows the comparison is fair:
event-hour trailing volatility (223 bp) is essentially identical to the
pool (225 bp) --- the effect is \emph{not} an artefact of posts arriving
in already-high-volatility hours.

The economic magnitude is modest --- roughly a 5--8\% relative increase
in realised volatility for one to four hours --- and we treat it as a
volatility/timing signal, not an alpha source.

\subsubsection{6.2 The content placebo --- the effect is specific to
market
content}\label{the-content-placebo-the-effect-is-specific-to-market-content}

To separate ``market content moves volatility'' from ``any post, or any
time-of-day, coincides with volatility,'' we run the identical pipeline
on the \emph{same author's non-market-relevant} posts --- 7,232 bursts,
five times the treatment count, so higher statistical power. Under the
primary inference these show \textbf{no} volatility effect: the 4-hour
event dummy is +1.04 bp (\emph{p} = 0.232) full-sample and +1.62 bp
(\emph{p} = 0.214) in 2025+, versus the treatment's +5.18 bp (\emph{p} =
0.007) and +6.33 bp (\emph{p} = 0.015) --- a roughly five-fold
difference in effect size, with the placebo indistinguishable from zero.
At 24 hours neither group shows a positive effect; the placebo's 24-hour
coefficient is in fact significantly \emph{negative} (−5.39 bp, \emph{p}
= 0.004) --- an unexplained large-sample artefact of the opposite sign
that we do not interpret as a post effect, and which mirrors the
treatment's own null 24-hour coefficient (−0.89 bp). The volatility
effect is therefore a property of the market \emph{content} of the posts
at the 1--4 hour horizon, not of the act of posting or of intraday
seasonality (which the null already matches).

\subsubsection{6.3 RQ2/RQ3 --- no directional predictability, no
tradeable
edge}\label{rq2rq3-no-directional-predictability-no-tradeable-edge}

The directional signal does \textbf{not} predict the sign of forward
returns: the continuous \texttt{signal} term is insignificant at every
asset and horizon (all \emph{p} \textgreater{} 0.4). No post-event
trading rule clears costs: across all twelve cells (three assets ×
\{all, policy-gated\} × \{LLM, FinBERT signals\}), the sign-of-signal
rule loses money net of the 0.24\% round trip, with out-of-sample
annualised Sharpe ratios between roughly −6 and −15. Critically, the
null holds under the \textbf{FinBERT-only} signal as well as the LLM
signal, so it is not an artefact of any LLM training-corpus leakage.
Trading every market-relevant burst is, mechanically, cost-prohibitive;
the absence of a directional edge is robust.

The Benjamini--Hochberg procedure flags eight ``surviving''
coefficients, but \textbf{all eight are the market-directive dummy} ---
i.e.~the \emph{n} = 2 directive class (below) --- and their large
magnitudes (+166 to +724 bp) and \emph{p} = 0.000 are the mechanical
artefact of a two-observation dummy landing on two large-move events,
not evidence of a generalisable effect. A negative one-hour coefficient
on the China/tariff topic (−8 to −17 bp, \emph{p} ≈ 0.02--0.03, not
surviving correction) is reported only as a consistency check: that
trade-war escalation is risk-off for crypto is already known, so we do
not report it as a finding.

\subsubsection{\texorpdfstring{6.4 The market-directive case study
(\emph{n} =
2)}{6.4 The market-directive case study (n = 2)}}\label{the-market-directive-case-study-n-2}

Two posts in the corpus are explicit market directives: ``I think it's
going very well\ldots The MARKETS are going to BOOM'' (3 April 2025) and
``THIS IS A GREAT TIME TO BUY!!! DJT'' (9 April 2025). Both preceded
large rallies; the 9 April post preceded the 90-day tariff-pause
announcement by roughly four hours, and Bitcoin rose
\textasciitilde5.6\% over the following four hours and
\textasciitilde4.0\% over 24 hours. We report these as a qualitative
illustration of the directive class and, as a documented matter of
public record, note that the 9 April post prompted public calls for
insider-trading and market-manipulation investigations. With \emph{n} =
2 we draw no statistical inference; the episode motivates the directive
taxonomy and the future collection of more such events, not a tradeable
claim.

\subsubsection{6.5 Macro panel}\label{macro-panel}

On the daily macro panel (December 2024 onward) we find no statistically
significant next-day effects of post activity or sentiment on any
series. The panel is short and daily, so this is low-power evidence
rather than evidence of absence.

\begin{center}\rule{0.5\linewidth}{0.5pt}\end{center}

\subsection{7. Discussion and
limitations}\label{discussion-and-limitations}

\textbf{What the evidence supports.} Posts with market content are
followed by a small, robust increase in 1--4 hour crypto volatility that
survives an endogenous-timing control and is absent for the same
author's non-market posts. They do not predict direction, and the
volatility effect is too small and too short-lived to trade after costs.
This is a coherent and, we think, credible picture: high-authority
communications inject short-lived uncertainty without providing a
directional edge to an outside reader.

\textbf{Methodological by-catch.} The content placebo also audited our
inference machinery: the day-block bootstrap null flagged a
\emph{negligible} +0.4 bp placebo ``effect'' as significant, revealing
that this null is anti-conservative at large \emph{n} (its standard
error shrinks faster than the dependence structure warrants). We
therefore treat the \textbf{cluster-robust regression as the primary
inference} and the bootstrap nulls as corroboration --- a choice the
placebo justifies empirically.

\textbf{Limitations (data ceilings, not fixable by method).} (i) The
three crypto assets co-move at \textasciitilde0.92, so they constitute
approximately one effect observed three times, not three independent
confirmations; we report Bitcoin as primary and the others as
robustness. (ii) Market-relevant bursts cluster heavily in the 2025--26
trade-policy regime, so external validity beyond that regime is
unestablished, and this bites the directional results hardest. (iii) The
market-directive class has \emph{n} = 2; it is a case study. (iv) Our
placebo is a \emph{content} (within-actor) control; a
\emph{different-actor} placebo --- another high-salience poster's
timestamps --- would further separate ``this actor's text'' from ``any
salient market commentator,'' but no clean negative-control feed was
readily available (a poster who also moves crypto would be a positive,
not negative, control). We log this as the priority extension. (v) The
LLM-derived labels are not bit-reproducible under the model alias used;
a local, pinned model would fix this and is planned.

\textbf{Future work.} The natural extensions are: a different-actor
placebo; a test of whether the volatility signal is monetisable through
an instrument matched to it (e.g.~options/variance products, absent from
our spot data); and --- once a 24/7 collector has grown the
directive/escalation event set well beyond \emph{n} = 2 --- a proper
test of whether such posts are predictable in advance from prevailing
conditions.

\begin{center}\rule{0.5\linewidth}{0.5pt}\end{center}

\subsection{8. Conclusion}\label{conclusion}

A sitting head of state's social-media posts, when they concern markets,
are followed by a modest, robust, content-specific increase in
short-horizon (1--4 hour) cryptocurrency volatility, but not by
predictable return direction, and not by a tradeable edge after costs.
The widely reported ``buy'' directives are a compelling but
statistically thin case study, not an edge. The result is a careful
\emph{measurement} rather than a strategy --- which is, given the data,
the honest place to stop.

\begin{center}\rule{0.5\linewidth}{0.5pt}\end{center}

\subsection{Reproducibility \&
disclosure}\label{reproducibility-disclosure}

All code, the pre-registration and its amendment, the verified citation
set, and a dated work-log of every decision are version-controlled in
the project repository; the frozen signal label set and pinned
post-archive snapshot are recorded with SHA-256 hashes in a signal
manifest. Analyses use a fixed random seed. The research was conducted
with substantial AI coding and analysis assistance; per standard journal
policy, AI systems are not listed as authors, and their role is
disclosed here and documented in full in the repository's work-log.
\emph{{[}Author names, affiliations, funding, and competing-interest
statements to be completed.{]}}

\subsection{References (working list --- see RELATED\_WORK\_VERIFIED.md
for the full annotated
set)}\label{references-working-list-see-related_work_verified.md-for-the-full-annotated-set}

\begin{itemize}
\tightlist
\item
  Born, J. A., Myers, D. H. \& Clark, W. J. (2017). Trump tweets and the
  efficient market hypothesis. \emph{Algorithmic Finance}, 6(3--4),
  103--109.
\item
  Burggraf, T., Fendel, R. \& Huynh, T. L. D. (2020). Political news and
  stock prices: evidence from Trump's trade war. \emph{Applied Economics
  Letters}, 27(18).
\item
  Gjerstad, P., Meyn, P. F., Molnár, P. \& Næss, T. D. (2021). Do
  President Trump's tweets affect financial markets? \emph{Decision
  Support Systems}, 147, 113577.
\item
  Klaus, J. \& Koser, C. (2021). Measuring Trump: the Volfefe Index and
  its impact on European financial markets. \emph{Finance Research
  Letters}, 38, 101447.
\item
  Krause, D. (2025). \emph{Trump-Era Cryptocurrency Policy Events and
  Bitcoin Returns}; \emph{Trump, Tokens, and Tailwinds}; \emph{Tariffs,
  Tokens, and Turmoil} (SSRN working papers).
\item
  Nishimura, Y., Dong, X. \& Sun, B. (2021). Trump's tweets: sentiment,
  stock market volatility, and jumps. \emph{Journal of Financial
  Research}, 44(3), 497--512.
\item
  Perico Ortiz, D. (2023). Economic policy statements, social media, and
  stock market uncertainty. \emph{Journal of Economics and Finance}, 47,
  333--367.
\end{itemize}

\end{document}
